\documentclass[12pt,a4paper]{article}
\usepackage[UTF8]{ctex}
\usepackage{amsmath,amssymb,amsthm}
\usepackage{geometry}

\geometry{margin=2.5cm}

\title{递归牛顿-欧拉算法后向递推公式推导}
\author{第8章补充说明}
\date{}

\begin{document}

\maketitle

\section{问题提出}

\textbf{问题}：在递归牛顿-欧拉算法中，后向递推公式是如何推导的？

\textbf{公式}：

\subsection{初始化（末端）}

\begin{equation}
f_{n+1} = 0, \quad \mu_{n+1} = 0
\end{equation}

\subsection{力和力矩递推}

对于每个关节$i = n, n-1, \ldots, 1$：
\begin{align}
f_i &= R_{i+1} f_{i+1} + m_i \dot{v}_{c,i} \tag{24} \\
\mu_i &= R_{i+1} \mu_{i+1} + r_i \times f_i + I_i \dot{\omega}_i + \omega_i \times (I_i \omega_i) \tag{25}
\end{align}

\subsection{关节力矩计算}

\begin{itemize}
\item 如果关节$i$是旋转关节：
\begin{equation}
\tau_i = \mu_i^T z_i \tag{26}
\end{equation}

\item 如果关节$i$是移动关节：
\begin{equation}
\tau_i = f_i^T z_i \tag{27}
\end{equation}
\end{itemize}

\section{符号说明}

\begin{itemize}
\item $f_i$：作用在连杆$i$上的力（在坐标系$\{i\}$中表示）
\item $\mu_i$：作用在连杆$i$上的力矩（在坐标系$\{i\}$中表示，相对于坐标系$\{i\}$的原点）
\item $f_{i+1}$：作用在连杆$i+1$上的力（在坐标系$\{i+1\}$中表示）
\item $\mu_{i+1}$：作用在连杆$i+1$上的力矩（在坐标系$\{i+1\}$中表示）
\item $R_{i+1}$：从坐标系$\{i+1\}$到坐标系$\{i\}$的旋转矩阵
\item $m_i$：连杆$i$的质量
\item $\dot{v}_{c,i}$：连杆$i$质心的线加速度（在坐标系$\{i\}$中表示）
\item $r_i$：从坐标系$\{i\}$原点到连杆$i$质心的向量（在坐标系$\{i\}$中表示）
\item $I_i$：连杆$i$在质心处的惯性张量（在坐标系$\{i\}$中表示）
\item $\omega_i$：连杆$i$的角速度（在坐标系$\{i\}$中表示）
\item $\dot{\omega}_i$：连杆$i$的角加速度（在坐标系$\{i\}$中表示）
\item $z_i$：关节$i$的旋转轴或移动方向（在坐标系$\{i\}$中表示，单位向量）
\item $\tau_i$：关节$i$的驱动力矩（旋转关节）或驱动力（移动关节）
\end{itemize}

\section{基础：刚体动力学}

\subsection{牛顿第二定律}

对于刚体，牛顿第二定律可以写成：
\begin{equation}
\mathbf{F} = m \mathbf{a}
\end{equation}

其中：
\begin{itemize}
\item $\mathbf{F}$：作用在刚体上的合力
\item $m$：刚体的质量
\item $\mathbf{a}$：刚体质心的加速度
\end{itemize}

\subsection{欧拉方程}

对于刚体的旋转运动，欧拉方程为：
\begin{equation}
\boldsymbol{\mu} = I \dot{\boldsymbol{\omega}} + \boldsymbol{\omega} \times (I \boldsymbol{\omega})
\end{equation}

其中：
\begin{itemize}
\item $\boldsymbol{\mu}$：作用在刚体上的合力矩（相对于质心）
\item $I$：刚体在质心处的惯性张量
\item $\dot{\boldsymbol{\omega}}$：刚体的角加速度
\item $\boldsymbol{\omega}$：刚体的角速度
\end{itemize}

\textbf{注意}：$\boldsymbol{\omega} \times (I \boldsymbol{\omega})$是科里奥利项，表示由于旋转产生的力矩。

\subsection{力的合成}

作用在连杆$i$上的总力等于：
\begin{equation}
f_i = f_{\text{传递}} + f_{\text{惯性}}
\end{equation}

其中：
\begin{itemize}
\item $f_{\text{传递}}$：从连杆$i+1$传递过来的力
\item $f_{\text{惯性}}$：连杆$i$的惯性力（$m_i \dot{v}_{c,i}$）
\end{itemize}

\subsection{力矩的合成}

作用在连杆$i$上的总力矩（相对于坐标系$\{i\}$的原点）等于：
\begin{equation}
\mu_i = \mu_{\text{传递}} + \mu_{\text{力臂}} + \mu_{\text{惯性}}
\end{equation}

其中：
\begin{itemize}
\item $\mu_{\text{传递}}$：从连杆$i+1$传递过来的力矩
\item $\mu_{\text{力臂}}$：由于力$f_i$不在原点产生的力矩（$r_i \times f_i$）
\item $\mu_{\text{惯性}}$：连杆$i$的惯性力矩（$I_i \dot{\omega}_i + \omega_i \times (I_i \omega_i)$）
\end{itemize}

\section{公式(24)：力的递推}

\subsection{物理意义}

作用在连杆$i$上的力由两部分组成：
\begin{enumerate}
\item 从连杆$i+1$传递过来的力：$R_{i+1} f_{i+1}$
\item 连杆$i$的惯性力：$m_i \dot{v}_{c,i}$
\end{enumerate}

\subsection{推导}

\textbf{步骤1：力的传递}

考虑连杆$i$和连杆$i+1$之间的连接。作用在连杆$i+1$上的力$f_{i+1}$（在坐标系$\{i+1\}$中表示）需要转换到坐标系$\{i\}$中。

\textbf{坐标变换}：

如果一个力在坐标系$\{i+1\}$中表示为$f_{i+1}$，那么在坐标系$\{i\}$中的表示为：
\begin{equation}
f_{\text{传递}} = R_{i+1} f_{i+1}
\end{equation}

其中$R_{i+1}$是从坐标系$\{i+1\}$到坐标系$\{i\}$的旋转矩阵。

\textbf{注意}：力是自由向量（free vector），不依赖于作用点，因此只需要旋转，不需要平移。

\textbf{步骤2：惯性力}

根据牛顿第二定律，连杆$i$的惯性力为：
\begin{equation}
f_{\text{惯性}} = m_i \dot{v}_{c,i}
\end{equation}

其中$\dot{v}_{c,i}$是连杆$i$质心的加速度（在前向递推中已计算）。

\textbf{步骤3：力的合成}

作用在连杆$i$上的总力为：
\begin{equation}
f_i = f_{\text{传递}} + f_{\text{惯性}} = R_{i+1} f_{i+1} + m_i \dot{v}_{c,i} \tag{24}
\end{equation}

\textbf{物理意义}：
\begin{itemize}
\item $R_{i+1} f_{i+1}$：从后续连杆传递过来的力
\item $m_i \dot{v}_{c,i}$：连杆$i$的惯性力（与加速度方向相反）
\end{itemize}

\section{公式(25)：力矩的递推}

\subsection{物理意义}

作用在连杆$i$上的力矩（相对于坐标系$\{i\}$的原点）由四部分组成：
\begin{enumerate}
\item 从连杆$i+1$传递过来的力矩：$R_{i+1} \mu_{i+1}$
\item 由于力$f_i$不在原点产生的力矩：$r_i \times f_i$
\item 由于角加速度产生的惯性力矩：$I_i \dot{\omega}_i$
\item 由于旋转产生的科里奥利力矩：$\omega_i \times (I_i \omega_i)$
\end{enumerate}

\subsection{推导}

\textbf{步骤1：力矩的传递}

作用在连杆$i+1$上的力矩$\mu_{i+1}$（在坐标系$\{i+1\}$中表示，相对于坐标系$\{i+1\}$的原点）需要转换到坐标系$\{i\}$中。

\textbf{坐标变换}：

如果一个力矩在坐标系$\{i+1\}$中表示为$\mu_{i+1}$，那么在坐标系$\{i\}$中的表示为：
\begin{equation}
\mu_{\text{传递}} = R_{i+1} \mu_{i+1}
\end{equation}

\textbf{注意}：力矩也是自由向量，只需要旋转，不需要平移。

\textbf{步骤2：力臂产生的力矩}

作用在连杆$i$上的力$f_i$作用在质心处（或通过质心），但我们需要计算相对于坐标系$\{i\}$原点的力矩。

\textbf{力矩定义}：

如果一个力$f$作用在点$P$，相对于点$O$的力矩为：
\begin{equation}
\mu = r_{OP} \times f
\end{equation}

其中$r_{OP}$是从点$O$到点$P$的向量。

在我们的情况下：
\begin{itemize}
\item 点$O$：坐标系$\{i\}$的原点
\item 点$P$：连杆$i$的质心
\item $r_i$：从坐标系$\{i\}$原点到质心的向量
\item $f_i$：作用在连杆$i$上的力
\end{itemize}

因此：
\begin{equation}
\mu_{\text{力臂}} = r_i \times f_i
\end{equation}

\textbf{步骤3：惯性力矩}

根据欧拉方程，连杆$i$的惯性力矩（相对于质心）为：
\begin{equation}
\mu_{\text{惯性,质心}} = I_i \dot{\omega}_i + \omega_i \times (I_i \omega_i)
\end{equation}

\textbf{关键问题}：这个力矩是相对于质心的，但我们需要相对于坐标系$\{i\}$原点的力矩。

\textbf{力矩的平移}：

如果惯性力矩$\mu_{\text{惯性,质心}}$是相对于质心的，那么相对于坐标系$\{i\}$原点的总惯性力矩为：
\begin{equation}
\mu_{\text{惯性}} = \mu_{\text{惯性,质心}} + r_i \times (m_i \dot{v}_{c,i})
\end{equation}

但是，我们已经将惯性力$m_i \dot{v}_{c,i}$包含在$f_i$中了，因此：
\begin{equation}
\mu_{\text{惯性}} = I_i \dot{\omega}_i + \omega_i \times (I_i \omega_i)
\end{equation}

\textbf{详细推导}：

考虑连杆$i$的质心$C_i$，其相对于坐标系$\{i\}$原点的位置为$r_i$。

\textbf{质心的加速度}：

质心的加速度为（在前向递推中已计算）：
\begin{equation}
\dot{v}_{c,i} = \dot{v}_i + \dot{\omega}_i \times r_i + \omega_i \times (\omega_i \times r_i)
\end{equation}

\textbf{质心的惯性力}：

根据牛顿第二定律，质心的惯性力为：
\begin{equation}
f_{\text{惯性,质心}} = m_i \dot{v}_{c,i}
\end{equation}

\textbf{质心的惯性力矩}：

根据欧拉方程，相对于质心的惯性力矩为：
\begin{equation}
\mu_{\text{惯性,质心}} = I_i \dot{\omega}_i + \omega_i \times (I_i \omega_i)
\end{equation}

\textbf{相对于原点的总力矩}：

相对于坐标系$\{i\}$原点的总力矩为：
\begin{align}
\mu_{\text{惯性}} &= \mu_{\text{惯性,质心}} + r_i \times f_{\text{惯性,质心}} \\
&= I_i \dot{\omega}_i + \omega_i \times (I_i \omega_i) + r_i \times (m_i \dot{v}_{c,i})
\end{align}

但是，由于$f_i$已经包含了$m_i \dot{v}_{c,i}$，而$r_i \times f_i$已经包含了$r_i \times (m_i \dot{v}_{c,i})$，因此我们只需要：
\begin{equation}
\mu_{\text{惯性}} = I_i \dot{\omega}_i + \omega_i \times (I_i \omega_i)
\end{equation}

\textbf{步骤4：力矩的合成}

作用在连杆$i$上的总力矩（相对于坐标系$\{i\}$的原点）为：
\begin{align}
\mu_i &= \mu_{\text{传递}} + \mu_{\text{力臂}} + \mu_{\text{惯性}} \\
&= R_{i+1} \mu_{i+1} + r_i \times f_i + I_i \dot{\omega}_i + \omega_i \times (I_i \omega_i) \tag{25}
\end{align}

\textbf{物理意义}：
\begin{itemize}
\item $R_{i+1} \mu_{i+1}$：从后续连杆传递过来的力矩
\item $r_i \times f_i$：由于力$f_i$不在原点产生的力矩
\item $I_i \dot{\omega}_i$：由于角加速度产生的惯性力矩
\item $\omega_i \times (I_i \omega_i)$：由于旋转产生的科里奥利力矩
\end{itemize}

\section{公式(26)和(27)：关节力矩计算}

\subsection{旋转关节（公式26）}

\textbf{物理意义}：

对于旋转关节，关节只能绕轴$z_i$旋转。因此，只有沿$z_i$方向的力矩分量需要由驱动器提供。

\textbf{推导}：

作用在连杆$i$上的总力矩$\mu_i$可以分解为：
\begin{equation}
\mu_i = \mu_{i,\parallel} + \mu_{i,\perp}
\end{equation}

其中：
\begin{itemize}
\item $\mu_{i,\parallel}$：沿$z_i$方向的分量（需要驱动器提供）
\item $\mu_{i,\perp}$：垂直于$z_i$的分量（由关节结构提供，不需要驱动器）
\end{itemize}

\textbf{投影}：

沿$z_i$方向的力矩分量为：
\begin{equation}
\mu_{i,\parallel} = (\mu_i^T z_i) z_i
\end{equation}

\textbf{关节力矩}：

关节$i$的驱动力矩为：
\begin{equation}
\tau_i = \mu_i^T z_i \tag{26}
\end{equation}

\textbf{物理意义}：
\begin{itemize}
\item 只有沿关节轴方向的力矩分量需要驱动器提供
\item 垂直于关节轴的力矩分量由关节结构（轴承、约束等）提供
\end{itemize}

\subsection{移动关节（公式27）}

\textbf{物理意义}：

对于移动关节，关节只能沿轴$z_i$移动。因此，只有沿$z_i$方向的力分量需要由驱动器提供。

\textbf{推导}：

作用在连杆$i$上的总力$f_i$可以分解为：
\begin{equation}
f_i = f_{i,\parallel} + f_{i,\perp}
\end{equation}

其中：
\begin{itemize}
\item $f_{i,\parallel}$：沿$z_i$方向的分量（需要驱动器提供）
\item $f_{i,\perp}$：垂直于$z_i$的分量（由关节结构提供，不需要驱动器）
\end{itemize}

\textbf{投影}：

沿$z_i$方向的力分量为：
\begin{equation}
f_{i,\parallel} = (f_i^T z_i) z_i
\end{equation}

\textbf{关节力}：

关节$i$的驱动力为：
\begin{equation}
\tau_i = f_i^T z_i \tag{27}
\end{equation}

\textbf{物理意义}：
\begin{itemize}
\item 只有沿关节轴方向的力分量需要驱动器提供
\item 垂直于关节轴的力分量由关节结构（导轨、约束等）提供
\end{itemize}

\section{初始化：末端条件}

\subsection{公式}

\begin{equation}
f_{n+1} = 0, \quad \mu_{n+1} = 0
\end{equation}

\subsection{物理意义}

\textbf{假设}：末端执行器没有受到外部力或力矩（或外部力/力矩已经在前向递推中考虑）。

\textbf{如果有外部力/力矩}：

如果末端执行器受到外部力$f_{\text{ext}}$和力矩$\mu_{\text{ext}}$（在末端执行器坐标系中表示），那么：
\begin{align}
f_{n+1} &= f_{\text{ext}} \\
\mu_{n+1} &= \mu_{\text{ext}}
\end{align}

\textbf{在递归算法中}：

通常，外部力/力矩可以：
\begin{enumerate}
\item 直接设置为$f_{n+1}$和$\mu_{n+1}$
\item 或者在前向递推中通过设置$\dot{v}_0$来考虑（对于重力）
\end{enumerate}

\section{完整算法流程}

\subsection{前向递推（从基座到末端）}

\begin{enumerate}
\item \textbf{初始化}：
\begin{align}
\omega_0 &= 0 \\
\dot{\omega}_0 &= 0 \\
v_0 &= 0 \\
\dot{v}_0 &= g
\end{align}

\item \textbf{对于每个关节$i = 1, \ldots, n$}：
\begin{itemize}
\item 计算$\omega_i$、$\dot{\omega}_i$、$v_i$、$\dot{v}_i$（根据关节类型）
\item 计算$\dot{v}_{c,i}$
\end{itemize}
\end{enumerate}

\subsection{后向递推（从末端到基座）}

\begin{enumerate}
\item \textbf{初始化}：
\begin{equation}
f_{n+1} = 0, \quad \mu_{n+1} = 0
\end{equation}

\item \textbf{对于每个关节$i = n, n-1, \ldots, 1$}：
\begin{itemize}
\item 计算$f_i$（公式24）
\item 计算$\mu_i$（公式25）
\item 计算$\tau_i$（公式26或27）
\end{itemize}
\end{enumerate}

\section{详细例子：单连杆机器人}

\subsection{场景}

\begin{itemize}
\item 基座固定在地面上
\item 一个旋转关节，连杆质量为$m$，长度为$L$
\item 连杆质心在连杆中心
\item 重力加速度：$g$
\end{itemize}

\subsection{前向递推}

\textbf{初始化}：
\begin{align}
\omega_0 &= 0 \\
\dot{\omega}_0 &= 0 \\
v_0 &= 0 \\
\dot{v}_0 &= g
\end{align}

\textbf{关节1（旋转关节）}：
\begin{align}
\omega_1 &= \dot{\theta}_1 z_1 \\
\dot{\omega}_1 &= \ddot{\theta}_1 z_1 \\
v_1 &= R_1^T (v_0 + \omega_0 \times p_1) = 0 \\
\dot{v}_1 &= R_1^T (\dot{v}_0 + \dot{\omega}_0 \times p_1 + \omega_0 \times (\omega_0 \times p_1)) = R_1^T g
\end{align}

\textbf{质心加速度}：
\begin{equation}
\dot{v}_{c,1} = \dot{v}_1 + \dot{\omega}_1 \times r_1 + \omega_1 \times (\omega_1 \times r_1)
\end{equation}

其中$r_1$是从关节1到质心的向量。

\subsection{后向递推}

\textbf{初始化}：
\begin{equation}
f_2 = 0, \quad \mu_2 = 0
\end{equation}

\textbf{关节1}：

\textbf{力}：
\begin{equation}
f_1 = R_2 f_2 + m_1 \dot{v}_{c,1} = m_1 \dot{v}_{c,1}
\end{equation}

\textbf{力矩}：
\begin{equation}
\mu_1 = R_2 \mu_2 + r_1 \times f_1 + I_1 \dot{\omega}_1 + \omega_1 \times (I_1 \omega_1)
\end{equation}

由于$R_2 \mu_2 = 0$，我们有：
\begin{equation}
\mu_1 = r_1 \times f_1 + I_1 \dot{\omega}_1 + \omega_1 \times (I_1 \omega_1)
\end{equation}

\textbf{关节力矩}：
\begin{equation}
\tau_1 = \mu_1^T z_1
\end{equation}

\section{总结}

\subsection{公式总结}

\textbf{后向递推}：
\begin{align}
f_i &= R_{i+1} f_{i+1} + m_i \dot{v}_{c,i} \tag{24} \\
\mu_i &= R_{i+1} \mu_{i+1} + r_i \times f_i + I_i \dot{\omega}_i + \omega_i \times (I_i \omega_i) \tag{25}
\end{align}

\textbf{关节力矩}：
\begin{itemize}
\item 旋转关节：$\tau_i = \mu_i^T z_i$ \tag{26}
\item 移动关节：$\tau_i = f_i^T z_i$ \tag{27}
\end{itemize}

\subsection{关键要点}

\begin{enumerate}
\item \textbf{力的传递}：从后续连杆传递过来的力需要坐标变换（旋转）
\item \textbf{惯性力}：根据牛顿第二定律，$f_{\text{惯性}} = m \dot{v}_c$
\item \textbf{力矩的传递}：从后续连杆传递过来的力矩需要坐标变换（旋转）
\item \textbf{力臂力矩}：由于力不在原点，产生力矩$r \times f$
\item \textbf{惯性力矩}：根据欧拉方程，包含角加速度项和科里奥利项
\item \textbf{关节力矩}：只有沿关节轴方向的分量需要驱动器提供
\item \textbf{递归结构}：从末端到基座，逐步计算每个连杆的力和力矩
\end{enumerate}

\subsection{计算复杂度}

\begin{itemize}
\item 前向递推：$O(n)$
\item 后向递推：$O(n)$
\item 总复杂度：$O(n)$（线性时间复杂度）
\end{itemize}

这比拉格朗日方法的$O(n^3)$要高效得多，特别适合实时控制应用。

\end{document}

